\section{Dimensions}

Why does physical space have three dimensions and not two or four. The substrate already forced four spacetime
dimensions from the rank of the octonion's rotation algebra, with one ordered direction for time and three for space, so
three is the count of unordered directions in the flat cusp. But there is a second, independent reason three is
selected, a dynamical one, and it is a genuine retrodiction.

The reason is orbits. A planet held to a star, an electron held to a nucleus, needs a stable closed orbit, an orbit that
comes back to itself rather than spiraling in or flying apart. Whether stable closed orbits exist at all depends on the
number of spatial dimensions. In three dimensions they exist, the familiar ellipses, because the inverse-square force is
exactly the one that closes. In two dimensions orbits precess away and never close, and in four or more dimensions they
are unstable, spiraling into the center or escaping to infinity at the slightest push. So only in three dimensions can
bound matter, atoms and planets and everything made of them, hold together at all.

An experiment confirms this directly, computing whether stable closed orbits exist as a function of the number of
dimensions, and finding that three is the only value that gives them, with two precessing and four and above unstable,
the control that gives the three its meaning. This is a retrodiction, not a free fit, because the answer could only have
been three if a world of bound matter is to exist, and the geometry independently delivered three from the rank of
$D_4$. Two roads, the rank of the rotation group and the stability of orbits, both land on three spatial dimensions, the
same way two roads landed on four spacetime dimensions and on the 24-cell. The mesh is selected to be the one place
where matter can bind.
